%\documentclass[11pt,a4paper,uplatex,dvipdfmx]{ujarticle} 		% for uplatex
\documentclass[11pt,a4j,dvipdfmx]{jarticle} 					% for platex
\input{pieces/form00_header} % pieces
\input{pieces/kakenhi7} % pieces
\input{pieces/form01_header} % pieces
\input{pieces/form02_2024_header} % pieces
\input{pieces/hook3} % pieces
%#Name: shingaku_fin
\input{pieces/form04_jsps_headers} % pieces
\input{pieces/form04_mext_headers} % pieces
\input{pieces/form04_shingaku_fin_header} % pieces
% ===== Global definitions for the Kakenhi form ======================
%　基本情報
%
%------ 研究課題名  -------------------------------------------
\newcommand{\研究課題名}{象の卵}

%----- 研究機関名と研究代表者の氏名-----------------------
\newcommand{\研究機関名}{逢坂大学}
\newcommand{\研究代表者氏名}{湯川秀樹}
\newcommand{\me}{\underline{\underline{H.~Yukawa}}} 
%---- 研究期間の最終年度 ----------------
\newcommand{\研究期間の最終元号年度}{6}  %令和で、半角数字のみ
%========================================

% user07_header
% ===== my favorite packages ====================================
% ここに、自分の使いたいパッケージを宣言して下さい。
\usepackage{wrapfig}
%\usepackage{amssymb}
%\usepackage{mb}
%\DeclareGraphicsRule{.tif}{png}{.png}{`convert #1 `dirname #1`/`basename #1 .tif`.png}
\usepackage{lineno}

% ===== my personal definitions ==================================
% ここに、自分のよく使う記号などを定義して下さい。
\newcommand{\klpionn}{K_L \to \pi^0 \nu \overline{\nu}}
\newcommand{\kppipnn}{K^+ \to \pi^+ \nu \overline{\nu}}

% ----- 業績リスト用 -------------
\newcommand{\paper}[6]{%
	% paper{title}{authors}{journal}{vol}{pages}{year}
	\item ``#1'', #2, #3 {\bf #4}, #5 (#6).			% お好みに合わせて変えてください。
}

\newcommand{\etal}{\textit{et al.\ }}
\newcommand{\ca}[1]{*#1}	% corresponding author;   \ca{\yukawa}  みたいにして使う
\newcommand{\invitedtalk}{招待講演}

\newcommand{\yukawa}{H.~Yukawa}					% no underline
%\newcommand{\yukawa}{\underline{\underline{H.~Yukawa}}}	% with 2 underlines
\newcommand{\tomonaga}{S.~Tomonaga}

\newcommand{\prl}{Phys.\ Rev.\ Lett.\ }		% よく使う雑誌も定義すると楽

% ===== 欄外メモ ==================
\newcommand{\memo}[1]{\marginpar{#1}}
%\renewcommand{\memo}[1]{}	% 全てのメモを表示させないようにするには、行頭の"%"を消す

\input{pieces/hook5} % pieces

\begin{document}
\input{pieces/hook7} % pieces
%#Split: 01_abstract  
%#PieceName: p01_abstract
\input{pieces/p01_abstract_00}
\section{計画概要}
%　　　　＜＜最大　1ページ＞＞

%s01_abstract_SG
%begin 研究目的の概要 ＝＝＝＝＝＝＝＝＝＝＝＝＝＝＝＝＝＝＝＝
\textbf{\\　　　　　＊＊＊　以下は、あくまで例です。真似しないでください。　＊＊＊\\
　　　　　＊＊＊　本文はもちろん、節の切り方や論理の組み方は　　　＊＊＊\\
　　　　　＊＊＊　ご自分の気に入ったスタイルで書いてください。　　＊＊＊}

	象の卵の研究の目的は．．．	
	\vspace*{10zw}	% （概要）と（本文）の間が10行程度になるよう、必要に応じて値を調整してください。
%end 研究目的の概要 ＝＝＝＝＝＝＝＝＝＝＝＝＝＝＝＝＝＝＝＝

\input{pieces/s01_abstract_SG_2nd_box} % pieces

%begin 研究計画の概要 ＝＝＝＝＝＝＝＝＝＝＝＝＝＝＝＝＝＝＝＝
	象の卵の研究計画と方法は．．．
	\vspace*{10zw}	% （概要）と（本文）の間が10行程度になるよう、必要に応じて値を調整してください。	
%end 研究計画の概要 ＝＝＝＝＝＝＝＝＝＝＝＝＝＝＝＝＝＝＝＝

\input{pieces/p01_abstract_01}

%#Split: 02_purpose  
%#PieceName: p02_purpose
\input{pieces/p02_purpose_00}
\section{研究成果取りまとめの目的}
%　　　　＜＜最大　1ページ＞＞

%s02_purpose_SGfin
%begin 研究目的 ＝＝＝＝＝＝＝＝＝＝＝＝＝＝＝＝＝＝＝＝
	 象の卵の研究目的は．．．

	\vspace{1cm}
	\begin{thebibliography}{99}
		\bibitem{teramura} 寺村輝夫、「ぼくは王様 - ぞうのたまごのたまごやき」.
	\end{thebibliography}
%end 研究目的 ＝＝＝＝＝＝＝＝＝＝＝＝＝＝＝＝＝＝＝＝

\input{pieces/p02_purpose_01}

%#Split: 03_past  
%#PieceName: p03_past
\input{pieces/p03_past_00}
\section{これまでの研究経過}
%　　　　＜＜最大　1ページ＞＞

%s03_past
%begin これまでの研究活動 ＝＝＝＝＝＝＝＝＝＝＝＝＝＝＝＝＝＝＝＝
	これまでは、シロナガスクジラの卵を探索してきたが、見つけられなかった。
%end これまでの研究活動 ＝＝＝＝＝＝＝＝＝＝＝＝＝＝＝＝＝＝＝＝

\input{pieces/p03_past_01}

%#Split: 04_review  
%#PieceName: p04_review
\input{pieces/p04_review_00}
\section{これまでの研究の評価}
%　　　　＜＜最大　1ページ＞＞

%s03_review
%begin これまでの研究の評価 ＝＝＝＝＝＝＝＝＝＝＝＝＝＝＝＝＝＝＝＝
	困難にもめげず、よくがんばったと思う。
%end これまでの研究の評価 ＝＝＝＝＝＝＝＝＝＝＝＝＝＝＝＝＝＝＝＝

\input{pieces/p04_review_01}

%#Split: 05_plan  
%#PieceName: p05_plan
\input{pieces/p05_plan_00}
\section{研究成果取りまとめの計画・方法}
%　　　　＜＜最大　1ページ＞＞

%s08_plan_1p.tex
%begin 研究計画 ＝＝＝＝＝＝＝＝＝＝＝＝＝＝＝＝＝＝＝＝
	象の卵の研究計画は．．．
%end 研究計画 ＝＝＝＝＝＝＝＝＝＝＝＝＝＝＝＝＝＝＝＝

\input{pieces/p05_plan_01}

%#Split: 06_publications  
%#PieceName: p06_publications
\input{pieces/p06_publications_00}
\section{研究成果の発表状況}
%　　　　＜＜最大　1ページ＞＞

% s14_pub
%begin 研究業績リスト ＝＝＝＝＝＝＝＝＝＝＝＝＝＝＝＝＝＝＝＝
	\begin{enumerate}
		\paper{Search for whale eggs}{\yukawa\ \etal}{Rev.\ Oceanic Mysteries}{888}{99}{2017}
			\label{pub:whale}
				
		\paper{Theory of Elephant Eggs}{\yukawa, Kara Juzo \etal}{\prl}{800}{800-804}{2005}
			\label{pub:theoegg}
				
		\paper{仔象は死んだ}{Kobo Abe}{安部公房全集}{26}{100-200}{2004}
		
		\paper{The Elephant's Child (象の鼻はなぜ長い)}{R.~Kipling}{Nature}{999}{777-799}{2003}

		\paper{You can't Lay an Egg If You're an Elephant}{F.~Ehrlich}
			{JofUR\\({\tt www.universalrejection.org})}{{\bf N/A}}{2002}
		
		% 下のように書いてもいいけど、めんどくさいし、表示の仕方を変えようとしたら大変。
		\item ``Egg of Elephant-Bird'', 
				\underline{A.~Cooper},
				Nature, {\bf 409}, 704-707 (2001).	% 	
		\input{jack_pub}	% << only for demonstration. Please delete it or comment it out.
	\end{enumerate}
%end 研究業績リスト ＝＝＝＝＝＝＝＝＝＝＝＝＝＝＝＝＝＝＝＝

\input{pieces/p06_publications_01}

%#Split: 04_rights  
%#PieceName: p04_rights
\input{pieces/p04_rights_00}
\section{人権の保護及び法令等の遵守への対応}
%　　　　＜＜最大　1ページ＞＞

% s09_rights
%begin 人権の保護及び法令等の遵守への対応 ＝＝＝＝＝＝＝＝＝＝＝＝＝＝＝＝＝＝＝＝
	象の卵のES細胞の培養、象のクローンの生成などは行わない。

	\LaTeX の便利な機能については、\texttt{egg\_***.tex} や\texttt{sample\_pdf/egg\_***.pdf}をご覧ください。
%end 人権の保護及び法令等の遵守への対応 ＝＝＝＝＝＝＝＝＝＝＝＝＝＝＝＝＝＝＝＝

\input{pieces/p04_rights_01}

%#Split: 99_tail
\input{pieces/hook9} % pieces
\end{document}

